\section{Results}\label{sec:results}
Unless noted, results are five-replicate medians on JUWELS; figures compare the matched-kernel synchronous baseline (labelled \emph{Synchronous}) against the proposed asynchronous method (\emph{Asynchronous}, ours). We lead with the computational evidence (RQ2), the paper's central claim, then verify that inference quality is not sacrificed (RQ1), and close with the realistic workload (RQ3).

\subsection{Computational advantage (RQ2)}\label{sec:results-hpc}
\paragraph{Straggler robustness.} This is the clearest illustration of the barrier's cost (Fig.~\ref{fig:straggler}). One worker is given a permanent post-evaluation delay of $0.1$\,s scaled by $0\times$ (control)$,1\times,5\times,10\times,20\times$. As the injected delay grows from the control to $20\times$, the synchronous baseline's throughput collapses by more than two orders of magnitude (median $\approx 8800 \to 60$ simulations/s) because each generation blocks on the straggler, whereas the asynchronous method routes work to available workers and holds essentially flat ($\approx 3900 \to 3800$ simulations/s); only at the $0\times$ control is the baseline faster, where a generation barrier is cheap for this near-instant simulator. Posterior quality is not paid for this throughput: measured against the analytic posterior, the asynchronous method's reported estimator sits at $0.008$--$0.011$ at every slowdown and the baseline's final-population estimator at $0.017$--$0.027$ (\S\ref{sec:results-validity}). We read that gap as a difference in the estimators' support --- a weighted full history against a $100$-particle final generation --- rather than as a difference between the samplers.

\paragraph{Isolating synchronization: a barrierized twin.} The comparison above is against a different algorithm, so it cannot by itself attribute the gain to the barrier (\S\ref{sec:limitations}, item vii). We therefore build a \emph{barrierized twin} of our own method: the identical propagator --- same proposals, same kernel weights, same archive rule, same reported estimator --- with a collective barrier inserted immediately before each breed, so that the asynchronous arm and the twin differ in \emph{nothing but} when a worker is permitted to proceed. Table~\ref{tab:twin} reports both on the straggler benchmark over five replicates. Only the twin needs a simulation-count budget: a collective barrier requires identical per-rank call counts, so a wall-clock stop --- which halts ranks independently, as the first one to reach the deadline --- deadlocks it. The asynchronous arm therefore runs wall-limited ($300$\,s of active wall-clock) and the twin runs to a fixed simulation count, and we compare the two as \emph{rates}: simulations per second of active wall-clock on both arms. The asynchronous arm holds $\approx3100$--$3900$ simulations/s at every slowdown, reproducing the throughput of the main experiment above; the twin runs at $48.7$ simulations/s already at the $0\times$ control --- the cost of synchronizing a $\approx4$\,ms simulator, where barrier latency alone dominates --- and degrades to $8.0$ as the straggler worsens, a factor of $402$ below the asynchronous arm at $20\times$. Barrier removal is therefore worth between one and two-and-a-half orders of magnitude of throughput on this workload, measured against our own algorithm rather than inferred from a cross-algorithm comparison.

\emph{And posterior quality?} It needs more care than the throughput does, because the statistic this comparison first used has no resolution on this benchmark. $W_1$ between the unweighted top-$k$ archive and a point mass at the true parameter is indeed $0.069$--$0.085$ on all three arms at every slowdown. But the analytic posterior here has standard deviation $\sigma_{\mathrm{obs}}/\sqrt{n_{\mathrm{obs}}} = 0.1$, so an \emph{exactly correct} posterior scores $0.090$ on that statistic: every arm is at or just below the score of a perfect answer, and the ones below it are there because a point-mass target rewards an over-concentrated archive. Agreement on that number is not evidence of anything, and we no longer report it. Evaluated instead on the estimator this paper actually reports --- the full-history AMIS posterior, against the analytic posterior --- the arms separate sharply (Table~\ref{tab:twin}). Through $5\times$ all three agree at $0.008$--$0.015$. At $10\times$ the fine twin's reported posterior moves to $0.74$ and at $20\times$ to $1.73$, with the coarse twin following one step behind at $1.07$; the asynchronous arm stays at $0.009$--$0.013$ throughout. The twins' \emph{archives} stay on target in every one of these runs --- it is the weights that degenerate, not the particles.

We do not read this as the barrier corrupting the estimator, because the two arms are not at a matched budget: at $20\times$ the twin completes $3{,}200$ evaluations against the asynchronous arm's $959{,}000$ in the same active wall-clock --- ten generations against roughly nine hundred --- and importance weights formed over ten generations are expected to be poor whatever produced them. What the data support is the weaker conjunction: \emph{at a fixed wall-clock budget} the barrier costs one to two-and-a-half orders of magnitude of throughput, and on this workload that deficit is large enough to cost the reported posterior as well. Separating the barrier's own effect on the weights from the effect of the evaluations it forgoes needs a matched-evaluation-count twin, which we have not run (\S\ref{sec:limitations}).

\emph{Is the twin's fine generation doing the work?} The twin above synchronizes every $W{=}16$ evaluations, one per worker --- a finer generation than the synchronous baseline's population of $100$ --- so one could suspect the ratio is an artifact of barrier \emph{frequency} rather than of synchronization as such. We tested this directly by re-running the twin with the barrier coarsened to one per $112$ evaluations ($7W$, matching the baseline's population), five replicates at every slowdown (Table~\ref{tab:twin}, third row). Granularity turns out to matter only where the straggler does not: at the $0\times$ and $1\times$ controls the coarse twin is $3.3\times$ and $2.3\times$ faster ($48.7\to162.8$ and $65.3\to148.4$ simulations/s), cutting the asynchronous lead from $80\times$ to $24\times$ and from $48\times$ to $21\times$; but at $5\times$, $10\times$ and $20\times$ the two granularities agree to three significant figures ($31.6$ vs $31.8$, $15.91$ vs $15.94$, $7.98$ vs $7.98$), leaving the ratios at $102$--$402\times$. The mechanism is straightforward: a barrier costs either its own latency or the wait for the slowest member of the generation, whichever is larger. Coarsening amortizes the latency over more evaluations, which helps when latency dominates, but every generation still has to wait for the straggler no matter how long the generation is --- so in the straggler regime the cost is granularity-independent. The advantage we report under a straggler is thus not an artifact of the twin's schedule, and the remaining caveat is one of workload rather than granularity: this is a deliberately adversarial near-instant simulator, and where per-evaluation cost amortizes the coordination (\S\ref{sec:results-cpm}) the barrier's relative cost is correspondingly smaller.

\paragraph{The same test under runtime heterogeneity.} A persistent straggler is the sharpest case but not the typical one, so we repeat the twin comparison on the heterogeneity benchmark, where every worker draws a fresh lognormal runtime multiplier at every evaluation rather than one rank being permanently slow. Both barrier granularities again, five replicates at each of the five heterogeneity levels, simulation-limited with the budget per level sized from measured rates (Table~\ref{tab:twin-hetero}). The picture differs from the straggler case in two instructive ways. First, the $\sigma{=}0$ control now costs essentially \emph{nothing}: the twin runs at $44.8$ simulations/s against the asynchronous arm's $47.6$, a factor of $1.1$, where on the straggler benchmark the same control already cost $24$--$80\times$. The difference is the simulator, not the barrier --- these evaluations take ${\approx}1$\,s rather than ${\approx}4$\,ms, so barrier latency is negligible against the work it separates. The barrier's cost on a cost-bearing simulator is therefore a pure \emph{variance} cost, and it appears only once there is variance to pay for: it grows monotonically to $10.1\times$ at $\sigma{=}1.5$ and $25.3\times$ at $\sigma{=}2$ at the baseline-matched granularity ($30.9\times$ at the finer one). Second, coarsening the barrier now \emph{does} help --- $1.2$ to $1.7\times$ across $\sigma\ge0.5$, against the $1.00\times$ it bought under the persistent straggler. That contrast is the mechanism made visible: coarsening trades a heavier per-barrier maximum against half as many barriers, which is a favourable trade when the slow draw is transient and a different worker each time, and no trade at all when the same rank is slow in every generation. Posterior quality on these runs needs the same care as on the straggler benchmark, and for a different reason. The archives agree: the twin's top-$k$ archive recovers the analytic posterior mean to within $0.003$--$0.034$ absolute error against the asynchronous arm's $0.007$--$0.011$ (the synchronous baseline's final population, by contrast, sits at $0.012$--$0.139$), and in $W_1$ to the analytic posterior the archives are at $0.014$--$0.019$ and $0.015$--$0.106$ respectively. But this study is simulation-limited at $2112$ evaluations, and at that budget the \emph{reported} full-history estimator has not converged on \emph{any} arm: its $W_1$ to the analytic posterior is $0.19$--$2.34$ across the four arms, against $2.42$ for the prior itself. We therefore read these runs as what they were built to be --- a throughput comparison --- and take no posterior-quality claim from them beyond the archives agreeing. The two vendored columns are named accordingly.

There is one further reading worth stating, because it runs opposite to the concern the twin was built to address. The worry about a cross-algorithm comparison is that it might \emph{overstate} the barrier's role. Here it understates it: the synchronous pyABC baseline's penalty grows only from $2.1\times$ to $4.8\times$ over the same heterogeneity range, so at $\sigma{=}2$ our own barrierized twin is $5.3\times$ slower than the external baseline at matched barrier granularity. The design difference that accounts for it is the one noted in \S\ref{sec:limitations}, item iv: pyABC closes a generation on the first $100$ evaluations to finish and discards the latecomers, so it never waits for the slowest draw, whereas the twin --- like our own propagator, minus asynchrony --- keeps every evaluation and therefore pays the full maximum. The barrier's unmitigated cost is thus larger than the measured gap over pyABC, and the headline comparison is conservative rather than flattering.

\begin{table}[t]
\caption{Barrierized twin on the straggler benchmark: median simulation throughput (simulations per second of active wall-clock) over five replicates. The twin is the \emph{same} propagator with a collective barrier before each breed, so the rows differ only in synchronization. Two barrier granularities are shown: one barrier per $W{=}16$ evaluations (one per worker, the finest possible generation) and one per $112$ evaluations ($7W$, matching the synchronous baseline's population of $100$). Coarsening the generation to the baseline's costs the twin nothing once the straggler binds --- the $5\times$/$10\times$/$20\times$ columns agree to three significant figures --- and buys a factor of $3.3$/$2.3$ only at the $0\times$/$1\times$ controls, where barrier latency rather than straggler waiting is the bottleneck. The lower block reports posterior quality as $W_1$ between each arm's \emph{reported} estimator (the full-history AMIS posterior) and the \emph{analytic} posterior. We do not report $W_1$ to a point mass at the true parameter here: the analytic posterior for this configuration has standard deviation $0.1$, so an exactly correct answer scores $0.090$ on that statistic and every arm scores $0.069$--$0.085$ --- it has no resolution on this benchmark, and it rewards an over-concentrated archive. Note that the arms are not at a matched evaluation budget (last block); see the discussion in \S\ref{sec:results-hpc}.}\label{tab:twin}
\centering
\begin{tabular}{@{}lccccc@{}}
\toprule
Straggler slowdown & $0\times$ & $1\times$ & $5\times$ & $10\times$ & $20\times$ \\
\midrule
Asynchronous (ours)              & 3888 & 3112 & 3241 & 3218 & 3205 \\
Barrierized twin, every $W{=}16$ & 48.7 & 65.3 & 31.6 & 15.9 & 8.0 \\
Barrierized twin, every $112$    & 162.8 & 148.4 & 31.8 & 15.9 & 8.0 \\
\midrule
Ratio, twin every $W$      & $80\times$ & $48\times$ & $103\times$ & $202\times$ & $402\times$ \\
Ratio, twin every $112$     & $24\times$ & $21\times$ & $102\times$ & $202\times$ & $402\times$ \\
\midrule
\multicolumn{6}{@{}l}{\emph{Reported posterior vs.\ analytic posterior, median $W_1$}} \\
Asynchronous (ours)              & 0.010 & 0.013 & 0.011 & 0.011 & 0.009 \\
Barrierized twin, every $W{=}16$ & 0.011 & 0.008 & 0.015 & 0.737 & 1.729 \\
Barrierized twin, every $112$    & 0.012 & 0.011 & 0.009 & 0.016 & 1.067 \\
\midrule
\multicolumn{6}{@{}l}{\emph{Evaluations completed (median); the arms are not budget-matched}} \\
Asynchronous (ours)              & 1166k & 934k & 972k & 965k & 959k \\
Barrierized twin (both)          & 75.1k & 44.5k & 12.3k & 6.3k & 3.2k \\
\botrule
\end{tabular}
\end{table}

\begin{table}[t]
\caption{Barrierized twin under \emph{runtime heterogeneity} (each worker draws a fresh lognormal runtime multiplier of scale $\sigma$ at every evaluation), $48$ workers, median throughput over five replicates. Same two barrier granularities as Table~\ref{tab:twin}: one per $W{=}48$ evaluations and one per $96$, the latter matching the synchronous baseline's population of $100$. Unlike the straggler benchmark, the $\sigma{=}0$ control costs almost nothing ($1.1\times$) because these evaluations take ${\approx}1$\,s and barrier latency is negligible against them --- the cost here is purely the wait for the slowest draw, and it grows with $\sigma$. Coarsening the barrier helps under transient heterogeneity ($1.2$--$1.7\times$), where under a persistent straggler it bought nothing. The synchronous pyABC baseline degrades far less than our own twin because it closes a generation on the first $100$ finishers and discards latecomers, so at $\sigma{=}2$ the twin is $5.3\times$ slower than the external baseline it stands in for. All rates are simulations per second of each arm's own elapsed wall-clock; at $\sigma{=}1.5$ and $2$ the baseline's runs end at ${\approx}54$ and ${\approx}52$\,s of the $60$\,s budget, so dividing by its actual elapsed time credits it with the higher rate and makes the comparison against it conservative.}\label{tab:twin-hetero}
\centering
\begin{tabular}{@{}lccccc@{}}
\toprule
Heterogeneity $\sigma$ & $0$ & $0.5$ & $1.0$ & $1.5$ & $2.0$ \\
\midrule
Asynchronous (ours)               & 47.6 & 42.1 & 28.9 & 17.6 & 11.4 \\
Barrierized twin, every $W{=}48$  & 44.8 & 15.1 & 4.37 & 1.06 & 0.37 \\
Barrierized twin, every $96$      & 45.1 & 19.8 & 6.50 & 1.75 & 0.45 \\
Synchronous pyABC baseline        & 23.1 & 18.9 & 9.91 & 4.35 & 2.39 \\
\midrule
Ratio, twin every $W$   & $1.1\times$ & $2.8\times$ & $6.6\times$ & $16.6\times$ & $30.9\times$ \\
Ratio, twin every $96$  & $1.1\times$ & $2.1\times$ & $4.4\times$ & $10.1\times$ & $25.3\times$ \\
Ratio, pyABC baseline   & $2.1\times$ & $2.2\times$ & $2.9\times$ & $4.1\times$ & $4.8\times$ \\
\botrule
\end{tabular}
\end{table}

\begin{figure}[t]
\centering
\includegraphics[width=0.6\linewidth]{figures/fig_straggler_throughput.pdf}
\caption{Throughput (median $+$ inter-quartile range over five replicates) versus straggler slowdown factor, including a $0\times$ no-straggler control. One worker carries a permanent post-evaluation delay of $0.1$\,s scaled by the factor. The synchronous baseline degrades sharply at the generation barrier (from ${\approx}8800$ sims/s at the control to ${\approx}60$ at $20\times$); the asynchronous method is robust (median ${\approx}2600$--$3900$ across the range, with no collapse). Note the logarithmic throughput axis.}\label{fig:straggler}
\end{figure}

\paragraph{Runtime heterogeneity.} Under stochastic log-normal runtime noise the asynchronous method sustains higher utilization (lower idle fraction) than the synchronous baseline, which idles at generation barriers (Fig.~\ref{fig:hetero}). Crucially, this utilization converts directly into inference quality: at a fixed wall-clock budget the synchronous baseline completes fewer and fewer simulations as the runtime spread $\sigma$ grows (it idles at every barrier) and its posterior-mean error to the analytic Gaussian posterior climbs several-fold (from ${\approx}0.05$ to ${\approx}0.28$), whereas the asynchronous method completes several-fold more simulations and holds its error essentially flat (${\approx}0.007$--$0.011$) (Fig.~\ref{fig:hetero-quality}). Removing the barrier therefore preserves not just the simulation rate but the \emph{inference} delivered per unit wall-clock.

\begin{figure}[t]
\centering
\includegraphics[width=0.6\linewidth]{figures/fig_hetero_idle.pdf}
\caption{Worker utilization loss under runtime heterogeneity: the fraction of worker wall-clock spent idle versus the runtime-noise spread $\sigma$ (five replicates; bands are inter-quartile ranges). The synchronous baseline idles at generation barriers (its idle fraction rises from ${\approx}0.51$ to ${\approx}0.79$ as heterogeneity grows) while the asynchronous method stays near-fully utilized (idle fraction ${\approx}0.01$ rising to ${\approx}0.27$).}\label{fig:hetero}
\end{figure}

\begin{figure}[t]
\centering
\includegraphics[width=\linewidth]{figures/fig_hetero_quality.pdf}
\caption{Inferential efficiency under runtime heterogeneity (Gaussian-mean benchmark, analytic posterior, fixed $60$\,s budget, five replicates; bands are inter-quartile ranges). Left: posterior-mean error to the analytic posterior versus runtime-noise spread $\sigma$; the synchronous baseline degrades sharply while the asynchronous method stays flat. Right: simulations completed within the budget; the synchronous baseline is starved by barrier idling as $\sigma$ grows, whereas the asynchronous method sustains several-fold higher completion.}\label{fig:hetero-quality}
\end{figure}

\paragraph{Parameter-coupled runtime (every-particle bias).} The study above slows simulations independently of their parameters; a sharper test couples runtime \emph{to the inferred parameter}, directly probing the every-particle bias of Limitation~(iv). On the Gaussian-mean benchmark ($48$ workers, $60$\,s budget, five replicates) we set the post-evaluation delay to $d(\mu)=\min\!\big(d_0\,e^{\,c(\mu-\mu_c)/\ell},\,d_{\max}\big)$ with base $d_0=0.02$\,s, centre $\mu_c=0$, length scale $\ell=0.15$, and cap $d_{\max}=0.5$\,s, and sweep the coupling strength $c\in\{0,0.5,1,2\}$ ($c=0$ recovers uniform runtime). Larger $c$ makes one region of parameter space evaluate far faster than its mirror image, so the asynchronous ``keep every accepted particle'' rule over-represents the fast region in the raw archive. The coupling has the intended effect: aggregate throughput falls as it steepens, by ${\approx}43\%$ for the asynchronous method from its peak (${\approx}1880/2010 \to 1150$ sims/s; it is flat between $c=0$ and $c=0.5$ and falls monotonically thereafter) and monotonically and more steeply for the barrier-bound baseline (${\approx}840 \to 250$; Fig.~\ref{fig:param-bias}a), confirming that the fast region dominates the raw sample counts. Yet the asynchronous posterior-mean error to the analytic Gaussian posterior does not increase as the coupling steepens; it stays ${\approx}0.01$ and comparable to the discard-latecomers synchronous baseline (Fig.~\ref{fig:param-bias}b). At the coupling levels tested the over-representation is measurable in the raw sampling rate but does not measurably distort the reported (unweighted top-$k$) posterior mean. A targeted ablation that removes the AMIS proposal adaptation ($S{=}0$) and re-runs the same coupling sweep leaves this robustness essentially unchanged --- the asynchronous posterior-mean error stays ${\approx}0.005$--$0.013$ with or without AMIS and shows no trend in the coupling --- which indicates that the raw archive is itself robust at these levels rather than that AMIS actively corrects the over-representation. Two finer tests that Limitation~(iv) motivates sharpen this rather than leave it open. First, we score the AMIS-\emph{reweighted} posterior --- the estimator the guarantees of \S\ref{sec:theory} concern, resampling the archive by the retroactive weight \eqref{eq:posterior-estimator} --- instead of the unweighted archive, and compare it against the no-AMIS variant ($S{=}0$) across the coupling sweep and a steeper regime (base delay $0.5$\,s, cap $8$\,s). The reweighted posterior recovers the analytic mean as accurately as the unweighted archive (absolute error ${\le}0.025$, no trend in the coupling) in both regimes and with or without AMIS; the reweighting's effect surfaces in the posterior \emph{spread}, not the mean --- it lowers the effective sample size (ESS fraction ${\approx}0.88$, vs ${\approx}0.99$ without AMIS), roughly triples the maximum normalised weight (${\approx}0.03$ vs ${\approx}0.01$), and widens the central $90\%$ interval by ${\approx}15\%$ --- so AMIS reshapes the tails but neither creates nor removes a mean bias at these coupling levels. Second, a drain-after-deadline variant lets every in-flight simulation finish past the wall-clock deadline and recomputes the posterior; comparing it against the censored run isolates the completion-time censoring that dividing by the proposal density does not correct. Draining moves the weighted posterior mean by at most $0.013$ --- comparable to the replicate-to-replicate spread (${\approx}0.01$) --- with no systematic direction, in both the mild and the steep regime. At these coupling levels the deadline censoring is therefore present in principle but negligible in effect, consistent with the sampler concentrating in the fast region near the posterior where few simulations are in flight at the deadline. We accordingly report the raw archive as robust here rather than claiming AMIS corrects the over-representation, and treat a formal censoring correction as future work (Limitation~(iv)).

\begin{figure}[t]
\centering
\includegraphics[width=0.92\linewidth]{figures/fig_param_bias.pdf}
\caption{Parameter-coupled runtime on the Gaussian-mean benchmark (analytic posterior; $48$ workers, $60$\,s budget, five replicates; bands are inter-quartile ranges). The post-evaluation delay grows with the inferred mean, so one region of parameter space is sampled far more often than its mirror image. (a) Aggregate throughput falls as the coupling steepens --- monotonically for the barrier-bound baseline, and from a peak at $c=0.5$ for the asynchronous method --- confirming that the fast region dominates the raw sample counts. (b) Posterior-mean error to the analytic posterior nonetheless does not increase with coupling and stays comparable between the asynchronous method and the discard-latecomers synchronous baseline, so the over-representation does not measurably distort the reported posterior at these coupling levels.}\label{fig:param-bias}
\end{figure}

\paragraph{Strong scaling.} Lotka--Volterra is deliberately an \emph{adversarial} simulator for a scaling study: each evaluation takes only $\approx 2$\,ms, so the simulator occupies only a small fraction of each worker's wall-clock and there is essentially no per-evaluation compute to distribute. A dedicated instrumented run measures this productive fraction directly (Fig.~\ref{fig:lv-timing}): at the $48$-worker throughput peak a worker spends only ${\approx}25\%$ of its time inside the simulator, leaving ${\approx}75\%$ for coordination and idle waiting, and that productive share falls to ${\approx}8\%$ and then ${\approx}2\%$ as the job spreads from one node to three to six. The mechanism is the method's single shared population: every completed evaluation is broadcast to all workers so that each can propose from the current archive, an all-to-all exchange whose per-arrival cost grows with the worker count. On a near-instantaneous simulator this exchange saturates and then actively slows the productive ranks. Aggregate throughput therefore rises almost linearly while the job fits on one node ($149\to369\to1983\to6240$ sims/s at $1\to4\to16\to48$ workers) and then \emph{declines} across fully packed multi-node jobs (worker counts that are exact multiples of $48$): $6240\to4195\to2615\to2088\to1729$ sims/s at $48\to144\to192\to240\to288$ workers, i.e.\ $1\to3\to4\to5\to6$ nodes, because adding ranks beyond one node lengthens the all-to-all exchange faster than it adds useful work. This is a property of coordinating a near-instantaneous simulator across a single shared population, not of asynchrony as such (precisely the regime in which asynchrony has least to offer), and it lifts once the simulator carries real per-evaluation cost that amortizes the coordination (the Cellular Potts result of \S\ref{sec:results-cpm}; Fig.~\ref{fig:scaling}, right). To keep the comparison fair at high core counts we give the synchronous ABC-SMC baseline a \emph{population size equal to the worker count} (a smaller population cannot occupy every core); it then uses all allocated cores and scales steadily, from ${\approx}440$ sims/s at $16$ workers to ${\approx}6000$ at $288$ (Table~\ref{tab:scaling}). On this free simulator the asynchronous method leads while the job stays within a few nodes (${\approx}4.5\times$ at $16$ workers, ${\approx}2\times$ at one full node), but its lead \emph{erodes} with scale and the fair synchronous baseline overtakes it beyond three nodes, where the growing all-to-all coordination cost dominates and the asynchronous method leaves most of the allocation idle. This is the honest boundary of the asynchronous advantage on a cheap, homogeneous simulator; it is largest exactly in the heterogeneous and straggler-prone regimes where barriers otherwise waste compute. The archive size sets where that boundary falls, because the per-arrival proposal reconstruction is $\mathcal{O}(k)$: enlarging the archive tenfold to $k=1000$ costs throughput at every worker count --- $39\%$ at the single-node peak ($6240\to3804$ sims/s) --- and brings the crossover with the fair baseline in from beyond three nodes to a single node (Table~\ref{tab:scaling}). Sweeping $k$ across $\{50,100,200,400,800,1000\}$ at every worker count resolves that cost into a smooth monotone decline in throughput and a single step in the crossover, which moves in from beyond three nodes to three nodes at $k\ge400$ (Table~\ref{tab:k-frontier}). The archive size is therefore a systems cost as well as a statistical knob --- but, put on the same axis as the calibration it buys, the two turn out \emph{not} to trade off over this range: the best-calibrated archive is also within $7\%$ of the fastest, and every archive larger than $k=100$ is both slower and worse calibrated (\S\ref{sec:results-validity}, Table~\ref{tab:k-frontier}).

\begin{table}[t]
\caption{Strong scaling on Lotka--Volterra: median simulation throughput (sims/s) at a 180\,s budget, over five replicates (Fig.~\ref{fig:scaling}, left, shows the inter-quartile spread). Asynchronous (ours) at both archive sizes vs.\ synchronous baseline. Multi-node points use \emph{fully packed} nodes (worker counts that are exact multiples of $48$); at counts above $100$ the synchronous baseline is given \emph{population size equal to the worker count} so it can occupy every core (a population of $100$ leaves the surplus idle), which is why it keeps scaling here. The asynchronous method peaks on one node and then declines as multi-node coordination dominates, so the baseline overtakes it beyond three nodes at $k=100$. Enlarging the archive tenfold costs throughput at every scale --- the per-arrival proposal reconstruction is $\mathcal{O}(k)$ --- and moves that crossover in to a single node.}\label{tab:scaling}
\centering
\begin{tabular}{@{}lcccccccc@{}}
\toprule
Workers & 1 & 4 & 16 & 48 & 144 & 192 & 240 & 288 \\
Nodes & $<$1 & $<$1 & $<$1 & 1 & 3 & 4 & 5 & 6 \\
\midrule
Synchronous & 105 & 96 & 437 & 2952 & 3757 & 4322 & 4539 & 6034 \\
Asynchronous ($k=100$) & 149 & 369 & 1983 & 6240 & 4195 & 2615 & 2088 & 1729 \\
Asynchronous ($k=1000$) & 106 & 133 & 620 & 3804 & 2804 & 1957 & 1587 & 1336 \\
\botrule
\end{tabular}
\end{table}

\begin{figure}[t]
\centering
\includegraphics[width=\linewidth]{figures/fig_scaling_combined.pdf}
\caption{Strong scaling: median simulation throughput versus worker count at a fixed wall-clock budget, asynchronous (ours) vs.\ a synchronous baseline given a population of $\max(100, W)$ so it can occupy every core: pyABC caps concurrency at the population size, so at $W\le100$ the default $k{=}100$ already exceeds the worker count and is fair as it stands, and above one full node the baseline is re-run at $k{=}W$. A ``$W$-worker'' point gives both arms $W$ cores, but the baseline's MPI sampler spends its root rank coordinating rather than simulating, so it has $W{-}1$ simulating ranks against the asynchronous method's $W$ --- at most a $2.1\%$ handicap at $W{=}48$, falling to $0.3\%$ at $W{=}384$, i.e.\ shrinking exactly where the reported gap grows. \textbf{Left:} Lotka--Volterra (near-instant simulator, $\approx 2$\,ms/eval; $180$\,s budget; categorical $x$-axis; bands are inter-quartile ranges over five replicates; node count $n$ in labels). With no per-evaluation compute to distribute, the asynchronous method's single-population all-to-all coordination cost grows with the worker count: throughput rises within one node, peaks there, and then declines across multi-node jobs, so its lead over the fair, steadily scaling baseline erodes and the baseline overtakes it beyond three nodes. \textbf{Right:} Cellular Potts (costly \texttt{cellsInSilico} simulator; $1800$\,s budget; log--log; medians over replicates). Real per-evaluation cost amortizes the coordination, so the asynchronous method tracks ideal-linear scaling (dotted) to eight nodes ($384$ workers, ${\approx}97\%$ parallel efficiency) while the barrier-bound baseline scales only sub-linearly, a ${\approx}5\times$ gap at $384$ workers that \emph{grows} with scale.}\label{fig:scaling}
\end{figure}

\begin{figure}[t]
\centering
\includegraphics[width=0.66\linewidth]{figures/fig_lv_timing.pdf}
\caption{Measured decomposition of asynchronous worker wall-time on Lotka--Volterra ($k=100$; $180$\,s budget) from a dedicated instrumented run. The productive simulation fraction (blue, annotated) --- the share of worker wall-clock spent inside the simulator --- collapses from $55\%$ at $16$ workers to $2\%$ at $256$, while coordination and idle (MPI exchange, per-arrival proposal reconstruction, serialization, waiting) dominate exactly as the job spreads beyond one node (16/48 = 1 node; 128 = 3 nodes; 256 = 6 nodes). The strong-scaling decline of Fig.~\ref{fig:scaling} is therefore a measured property of coordinating a near-instantaneous simulator, not an inferred one.}\label{fig:lv-timing}
\end{figure}

\subsection{Posterior validity (RQ1)}\label{sec:results-validity}
Having established the systems advantage, we ask what it costs in inference quality. Across calibration and matched-budget posterior recovery the reported streaming posterior --- the full-history AMIS estimator \eqref{eq:posterior-estimator} --- is comparable to the matched synchronous baseline in low dimensions and under multimodality, and in four dimensions errs toward conservatism (it mildly over-covers) rather than over-confidence.

\paragraph{Calibration.} Table~\ref{tab:sbc} reports SBC empirical coverage on the Gaussian-mean model (1000 trials, matched simulation budget), for the reported full-history AMIS posterior. The asynchronous method tracks the nominal level closely ($0.51/0.80/0.90/0.94$ at the $0.50/0.80/0.90/0.95$ levels), within Monte Carlo error of nominal at every level; the synchronous baseline under-covers at every level ($0.35/0.63/0.75/0.83$), i.e.\ it is over-confident. At 1000 trials a coverage estimate carries a Monte Carlo standard error of only ${\approx}\,\pm0.007$--$0.016$, well below the ${\approx}0.12$--$0.17$ gap between the methods, so the ordering is resolvable rather than noise: on this benchmark the streaming estimator is \emph{better} calibrated than the matched synchronous baseline. Reporting the estimator over the full evaluated history rather than the top-$k$ archive alone matters even here: in a paired $500$-trial run that changes only the reported support, the truncated archive mildly under-covers at the top levels ($0.89/0.92$ at the $0.90/0.95$ levels), which the full-history report corrects to nominal --- a reporting-support effect that grows with dimension, as the multimodal and four-dimensional tests below make explicit. This test is on a one-dimensional model with an analytic posterior; the higher-dimensional and multimodal settings where the archive/proposal approximations are most likely to bite are probed separately below. Figure~\ref{fig:sbc} shows the corresponding coverage curve and rank histogram.

\paragraph{Calibration under multimodality and in higher dimensions.} To probe exactly those settings, we run SBC (1000 trials, matched simulation budget) on two harder targets, reporting the same full-history AMIS posterior. The first is a \emph{multimodal} model $y_i\sim\mathcal N(\theta^2,\sigma^2)$: the sign ambiguity gives a symmetric bimodal posterior at $\pm\sqrt{\bar y}$, the case a top-$k$ archive is most at risk of collapsing onto a single mode. It does not collapse --- the archive retains \emph{both} modes in $85.7\%$ of trials and the mode containing the true value in $86.7\%$ (the remainder are largely near-$\theta{=}0$ trials where the two modes merge), the importance weights stay near-uniform (median effective-sample-size fraction $0.99$, median maximum normalised weight $0.011$), and coverage tracks the nominal level ($0.50/0.82/0.91/0.96$), calibrated within Monte Carlo error. The second is the \emph{four-dimensional} g-and-k model, where the reporting-support choice bites. The reported full-history posterior is \emph{conservative}: it mildly over-covers at every level and parameter (e.g.\ empirical coverage $0.62/0.91/0.96/0.99$ at the $0.50/0.80/0.90/0.95$ levels for the location parameter $A$, and similarly for $B,g,k$), because the finite-$\epsilon$ smooth-ABC posterior it reports is over-dispersed at the loose bandwidth a four-dimensional budget reaches. The efficient top-$k$ archive used for online proposal sits at the opposite end: it is over-concentrated and would instead \emph{under}-cover ($0.85/0.89$ at the $0.90/0.95$ levels for $A$), with healthy weights throughout (median effective-sample-size fraction $0.81$), so the miscalibration is a reporting-support effect, not weight degeneracy. The calibrated posterior is therefore \emph{bracketed} between the over-concentrated archive core and the over-dispersed full history, and the amount of history reported is the dispersion knob: widening the reported support from the archive toward the full history raises coverage monotonically and continuously through nominal --- for $A$, $0.85/0.89 \to 0.93/0.96 \to 0.95/0.98 \to 0.97/0.99$ at the $0.90/0.95$ levels for reported sizes $M=200/500/1500/\text{full}$, crossing nominal near $M\approx400$--$800$ (the shape parameters $B$ and $k$ are already near-nominal at $M{=}500$, e.g.\ $0.89/0.94$ for $B$). This is a property of \emph{how much of the history is reported}, not of the AMIS weighting (whose importance weights stay well-conditioned) or of the budget: doubling the simulation budget to $80{,}000$ per trial leaves the archive coverage essentially unchanged ($0.85/0.89$ at the $0.90/0.95$ levels for $A$, within Monte Carlo error for $B,g,k$). We accordingly report the streaming estimator, evaluated over the full evaluated history, as well calibrated in one dimension and under multimodality, and conservative --- over-covering rather than over-confident --- in four dimensions; an intermediate reported support ($M\approx400$--$800$) or a larger archive tighten the four-dimensional report toward exactly nominal.

\paragraph{How calibration depends on the archive size and the snapshot buffer.} The tests above fix $k$ and $S$ per benchmark, which leaves open whether the two efficiency knobs must be retuned as the dimension grows. To settle this we sweep them directly on a $d$-dimensional Gaussian-mean target, $y_i\sim\mathcal N(\theta,\sigma^2 I_d)$ with $\theta\in\mathbb R^d$ under a flat box prior, whose posterior is known exactly at every $d$ (SBC against that exact posterior returns $0.496/0.794/0.905/0.952$, confirming the target itself is calibrated, so any departure below is a property of the method). The grid is $d\in\{1,2,4,8,16\}\times k\in\{50,100,200,400,800\}\times S\in\{0,5,20,50\}$, $1000$ trials per cell, reporting the same full-history estimator (Table~\ref{tab:ks-sweep}).

The two knobs behave quite differently. The snapshot buffer $S$ has a threshold and then saturates: at $S=0$ the estimator is over-confident in low dimensions (mean deviation $-0.088$ at $d=1,2$), $S=5$ removes essentially all of it ($-0.001$), and $S=20$ and $S=50$ add nothing further --- at any dimension. The value of a richer balance-heuristic denominator therefore does \emph{not} grow with $d$ over this range, and the fixed $S=20$ we use throughout sits comfortably above the knee. The archive size $k$ is the opposite: its effect is large, strongly dimension-dependent, and \emph{not} monotone. Small archives can be severely over-confident in intermediate dimensions --- at $d=8,k=50$ empirical coverage is $0.29/0.52/0.63/0.71$ against the nominal $0.50/0.80/0.90/0.95$, a failure that persists at every $S$ --- while intermediate archives over-cover in the same regime ($+0.18$ to $+0.20$ at $k=200$--$400$). We note that this does not follow the floor a naive count of covariance parameters would predict: the perturbation covariance is a full $d\times d$ matrix estimated from the archive, suggesting $k\gtrsim d(d{+}1)/2$, yet $d=8$ fails at $k=50>36$ while $d=16$ remains within $0.06$ of nominal at $k=50<136$. The reading this supports is that $d=16$ is protected by its own difficulty: with $20{,}000$ simulations over a $[-5,5]^{16}$ prior the sampler barely concentrates, so the reported posterior stays broad and covers almost by default, whereas $d=8$ has enough signal to concentrate but not enough to concentrate correctly. Two further runs test that reading rather than leaving it as narrative. First, it predicts that a \emph{higher} dimension should look benign again: at $d=32$ the coverage is $0.55$--$0.57$, $0.83$--$0.84$, $0.92$, $0.96$ at the four levels for every $k\in\{50,100,800\}$ --- mildly conservative and nowhere near the $d=8$ failure, even at $k=50$. Second, the failure is not a budget shortfall: doubling the simulation budget at $d=8,k=50$ makes coverage slightly \emph{worse} ($0.26/0.47/0.58/0.67$ against $0.29/0.52/0.63/0.71$ at the original budget), exactly as expected if the extra simulations sharpen an already mislocated posterior rather than repairing it. We therefore do not offer a closed-form rule for $k$; the practical guidance is that the archive size is the dispersion knob (\S\ref{sec:results-validity}), that its safe range is problem- and dimension-specific, and that it should be chosen by checking calibration directly rather than by a formula.

\paragraph{The archive size does not, in this range, trade calibration against throughput.} Because $k$ has both a statistical effect (above) and a systems cost --- per-arrival proposal reconstruction is $\mathcal{O}(k)$, so a larger archive is slower at every worker count (\S\ref{sec:results-hpc}) --- it is natural to expect the two to pull against each other, and to have to buy calibration with throughput. Measuring both on one axis shows that expectation is wrong over the range we can test. Table~\ref{tab:k-frontier} puts the throughput cost of $k$ next to its worst-case miscalibration over the dimensions swept, using a throughput grid matched on everything but $k$ (Lotka--Volterra, $180$\,s wall cap, five replicates, worker counts $1$--$288$: $k\in\{100,1000\}$ from the main scaling run and $k\in\{50,200,400,800\}$ from a dedicated $k$-frontier run sharing that configuration). Of the six archive sizes, only two are Pareto-optimal, and they are the two smallest: $k=100$ is simultaneously the best-calibrated setting tested (worst deviation $0.064$, at $d=4$) and within $7\%$ of the fastest ($6240$ vs $6684$ simulations/s at the single-node peak). Every larger archive is dominated by it on \emph{both} axes at once --- $k=200$ is $5\%$ slower and nearly three times worse calibrated ($0.180$), $k=400$ is $14\%$ slower and $0.195$, and $k=800$ costs $31\%$ of throughput to end up still slightly worse calibrated than $k=100$ ($0.074$). The only genuinely faster archive, $k=50$, is the catastrophic one ($0.250$ at $d=8$). So the practical shape of the $k$ decision is not a frontier to negotiate but a floor to clear: go far enough above the small-archive failure region, and then stop, because further enlargement buys no calibration and costs throughput monotonically.

Two qualifications keep this from being read too broadly. The join is on $k$ alone and the two axes come from different benchmarks by design --- calibration needs a target whose posterior is known at every $d$, throughput needs the regime where an $\mathcal{O}(k)$ per-arrival cost is most exposed. That second choice makes the throughput column an \emph{upper} bound on the relative cost of $k$: Lotka--Volterra evaluations take ${\approx}2$\,ms, so coordination dominates there, and on a cost-bearing simulator the per-evaluation work amortizes the archive operation (\S\ref{sec:results-cpm}). The conclusion is therefore conservative in the right direction --- even where $k$ is most expensive, the best-calibrated archive is nearly free. What $k$ \emph{does} change at scale is where the fair synchronous baseline overtakes the asynchronous method, and there it acts as a step rather than a smooth trade: the crossover sits beyond three fully packed nodes for $k\le200$ and moves in to three nodes for $k\ge400$ (last column).

\begin{table}[t]
\caption{Calibration sensitivity to the archive size $k$ and the AMIS snapshot buffer $S$ on the $d$-dimensional Gaussian-mean target ($1000$ trials per cell; entries are the mean signed deviation of empirical from nominal coverage over the $0.50/0.80/0.90/0.95$ levels and the $d$ parameters). Negative is over-confident, positive is conservative; $0.000$ is exactly calibrated. \textbf{Left:} varying $k$ at $S=20$. \textbf{Right:} varying $S$ at $k=100$.}\label{tab:ks-sweep}
\centering
\begin{tabular}{@{}lccccc@{}}
\toprule
& \multicolumn{5}{c@{}}{archive size $k$ (at $S=20$)} \\
\cmidrule(l){2-6}
& 50 & 100 & 200 & 400 & 800 \\
\midrule
$d=1$  & -0.014 & -0.003 & +0.006 & +0.006 & +0.074 \\
$d=2$  & -0.020 & -0.006 & +0.009 & +0.030 & +0.072 \\
$d=4$  & +0.025 & +0.064 & +0.094 & +0.133 & +0.063 \\
$d=8$  & \textbf{-0.250} & +0.055 & +0.180 & +0.195 & +0.061 \\
$d=16$ & -0.047 & +0.028 & +0.049 & +0.053 & +0.055 \\
\botrule
\end{tabular}
\vspace{0.6em}

\begin{tabular}{@{}lcccc@{}}
\toprule
& \multicolumn{4}{c@{}}{snapshot buffer $S$ (at $k=100$)} \\
\cmidrule(l){2-5}
& 0 & 5 & 20 & 50 \\
\midrule
$d=1$  & \textbf{-0.088} & -0.001 & -0.003 & -0.005 \\
$d=2$  & \textbf{-0.088} & -0.001 & -0.006 & -0.003 \\
$d=4$  & +0.039 & +0.068 & +0.064 & +0.064 \\
$d=8$  & +0.065 & +0.006 & +0.055 & +0.074 \\
$d=16$ & -0.010 & +0.010 & +0.028 & +0.034 \\
\botrule
\end{tabular}
\end{table}

\begin{table}[t]
\caption{The archive size on both of its axes at once. \textbf{Throughput}: median asynchronous throughput at the single-node peak ($48$ workers) on Lotka--Volterra, $180$\,s wall cap, five replicates ($k\in\{100,1000\}$ from the main scaling run, $k\in\{50,200,400,800\}$ from the matched $k$-frontier run), and its change relative to the $k{=}100$ we use throughout. \textbf{Calibration}: worst absolute deviation of empirical from nominal coverage over the dimensions of the $d$-dimensional Gaussian-mean sweep at $S{=}20$ (Table~\ref{tab:ks-sweep}), with the dimension attaining it. \textbf{Crossover}: fully packed nodes at which the fair synchronous baseline (population $=$ worker count) overtakes the asynchronous arm. Only $k{=}50$ and $k{=}100$ are Pareto-optimal; every larger archive is both slower \emph{and} worse calibrated than $k{=}100$. The sweep did not include $k{=}1000$.}\label{tab:k-frontier}
\centering
\begin{tabular}{@{}lccccc@{}}
\toprule
$k$ & sims/s at $W{=}48$ & vs.\ $k{=}100$ & worst $|\Delta\text{cov}|$ & at & baseline overtakes \\
\midrule
50   & 6684 & $+7\%$  & \textbf{0.250} & $d{=}8$ & beyond 3 nodes \\
100  & 6240 & ---     & \textbf{0.064} & $d{=}4$ & beyond 3 nodes \\
200  & 5904 & $-5\%$  & 0.180 & $d{=}8$ & beyond 3 nodes \\
400  & 5393 & $-14\%$ & 0.195 & $d{=}8$ & 3 nodes \\
800  & 4331 & $-31\%$ & 0.074 & $d{=}1$ & 3 nodes \\
1000 & 3804 & $-39\%$ & ---   & ---     & 3 nodes \\
\botrule
\end{tabular}
\end{table}

\begin{table}[t]
\caption{SBC empirical coverage on Gaussian-mean inference (parameter $\mu$, 1000 trials; Monte Carlo standard error ${\approx}\,\pm0.007$--$0.016$), for the reported full-history AMIS posterior. Closer to the nominal level is better.}\label{tab:sbc}
\centering
\begin{tabular}{@{}lcccc@{}}
\toprule
Method & 0.50 & 0.80 & 0.90 & 0.95 \\
\midrule
Synchronous baseline & 0.35 & 0.63 & 0.75 & 0.83 \\
Asynchronous (ours) & 0.51 & 0.80 & 0.90 & 0.94 \\
\botrule
\end{tabular}
\end{table}

\begin{figure}[t]
\centering
\includegraphics[width=0.48\linewidth]{figures/fig_sbc_coverage.pdf}\hfill
\includegraphics[width=0.48\linewidth]{figures/fig_sbc_rank.pdf}
\caption{Simulation-based calibration on Gaussian-mean inference (1000 trials; the reported full-history AMIS posterior): empirical coverage versus nominal level (left) and rank histogram (right). The asynchronous estimator is close to calibrated.}\label{fig:sbc}
\end{figure}

\paragraph{Archive concentration.} We first report the top-$k$ concentration diagnostic, on all three benchmarks and for completeness with the earlier literature, and then turn to the estimator the paper reports. Read as what it is --- distance from the archive to a point mass, which rewards concentration whether or not the posterior is right (\S\ref{sec:experiments-suite}) --- it says the following. On g-and-k the asynchronous method concentrates slightly \emph{tighter} than the synchronous baseline at matched wall-clock budget ($0.078$ vs.\ $0.090$, inter-quartile ranges $[0.077,0.083]$ and $[0.082,0.093]$), and on Lotka--Volterra the two are comparable with the asynchronous method marginally tighter ($0.316$ vs.\ $0.332$, overlapping inter-quartile ranges); on Lotka--Volterra both numbers rest on the ${\approx}2\%$ of simulations that do not go extinct (\S\ref{sec:experiments-suite}). On the expensive Cellular Potts simulator the ordering reverses and the asynchronous method is \emph{modestly but consistently behind}: the async$-$sync difference is positive across $87.5\%$ of the budget, ending at ${+}0.043$ with an inter-quartile envelope of $[0.019,0.052]$ that excludes zero, and it never exceeds ${\approx}0.06$ against a Wasserstein scale of ${\approx}0.4$ (Figs.~\ref{fig:posterior-recovery}c and~\ref{fig:cpm-diff}). One property of this benchmark bounds how that gap should be read: the Wasserstein-to-reference \emph{rises} with compute for both methods (asynchronous $0.37\to0.48$), because the Cellular Potts parameters are only weakly identified (\S\ref{sec:results-cpm}) and the comparison is against a reference point rather than a known posterior, so a diffuse cloud scores better than a posterior that has concentrated slightly off that reference. The metric therefore penalizes concentration on this target, and the asynchronous method --- which completes ${\approx}2\times$ the evaluations in the same budget --- concentrates fastest. Per-parameter posteriors and joint marginals are consistent with the known posteriors on the benchmarks that have one, and on Cellular Potts they are consistent for one parameter of two: the motility reference sits inside the central $90\%$ of both methods' marginals (at their $87$th and $91$st percentiles), whereas the division-rate reference falls below \emph{every} one of the $500$ reported draws for the asynchronous method and for the synchronous baseline alike (5th percentiles $0.19$ and $0.21$ against a reference of $0.05$). Both methods therefore exclude the division-rate reference outright rather than merely failing to resolve it, and rejection ABC places it only at its own $5$th percentile while remaining nearly uniform. This is the identifiability property of the benchmark asserting itself and it is common to all three methods (\S\ref{sec:results-cpm}); it is the reason no quality claim on Cellular Potts should be read as a bound on attainable accuracy. On the Gaussian-mean benchmark the asynchronous method recovers the posterior mean to within $\approx0.002$--$0.038$ absolute error, comparable to the synchronous baseline's $\approx0.006$--$0.021$ and far tighter than a rejection-ABC reference; the SBC calibration above is unaffected.

\paragraph{Scoring the reported estimator against a reference posterior.} The curves above measure the top-$k$ archive against a point mass at the true parameter, which is a concentration diagnostic and not posterior quality (\S\ref{sec:experiments-suite}). On the two benchmarks that admit a reference posterior we can do better and score the estimator this paper actually reports --- the retroactive AMIS posterior of \eqref{eq:posterior-estimator}, the object Theorem~\ref{thm:consistency} is about --- against that reference. We do it over time by replaying \texttt{extract\_posterior} on the prefix of the history that had completed by each checkpoint, which is what the run would have reported had it stopped there; the synchronous baseline is scored on its last completed generation and rejection ABC on its accepted set, so each method is judged on its own reported object (Fig.~\ref{fig:reported-recovery}). The reference is the analytic posterior for Gaussian-mean, and for g-and-k the posterior given the same seven octile summaries the discrepancy is built from --- \emph{not} the full-data posterior, which would measure the summaries' information loss rather than the sampler (Appendix~\ref{app:protocol}).

Two things come out of it, and the second was not what we went looking for.

\emph{First, the estimator recovers the posterior in one dimension and is beaten in four.} On Gaussian-mean the reported AMIS posterior sits at $W_1 = 0.008$--$0.016$ from the analytic posterior across the $300$\,s budget, against the synchronous baseline's $0.019$--$0.028$ and rejection ABC's $2.35$: the streaming estimator is about twice as accurate as the baseline's final population at every checkpoint. On g-and-k the ordering reverses and the margin is similar in size the other way. The baseline holds $0.021$--$0.043$ against the reference posterior over the $600$\,s budget and ends at $0.029$ (inter-quartile range $[0.021,0.029]$); the asynchronous method ranges over $0.042$--$0.296$ and ends at $0.069$ ($[0.037,0.073]$), with both far below rejection ABC's $1.01$. We report this plainly because the concentration diagnostic says the opposite on the same benchmark --- there the asynchronous archive is the tighter of the two ($0.078$ vs.\ $0.090$) --- and the disagreement is the point of measuring the reported estimator rather than the archive.

This is also a claim the point-mass metric was incapable of supporting in either direction: on the Gaussian configuration it floors at $0.090$, which is what an \emph{exactly correct} posterior scores, and every method sat at or below it (\S\ref{sec:results-hpc}).

\emph{Second, the effective sample size does not grow with the budget.} The reported posterior's ESS is $237$--$297$ particles at every checkpoint of the Gaussian run while the evaluated history grows from $1.0\times10^5$ to $1.3\times10^6$; the run's own stored weights give $261$--$303$ across the five replicates. Thirteen times the compute buys no additional effective particles, which is also why the accuracy curve above is flat rather than improving. The mechanism is not asynchrony --- it reproduces in a serial driver --- and it is not a defect of the replay. It is the interaction between where the estimator is reported and how the bandwidth is chosen: the posterior is reported at the tightest bandwidth the run reached, that bandwidth is set by ESS-retention bisection against the archive and therefore tracks the archive downwards, and so the set of particles carrying non-negligible kernel weight stays at a few multiples of $k$ while everything outside it is weighted to near zero. A controlled sweep makes the scaling explicit: over $k \in \{50,100,200\}$ and $n \in \{5\times10^3, 2\times10^4, 8\times10^4\}$, $\mathrm{ESS}/k$ stays within $1.9$--$3.6$ at all nine points while $\mathrm{ESS}/n$ falls like $1/n$. The full history enters the mixture denominator; it does not, in any effective sense, enter the numerator.

Dimension changes the character of this, and it is what explains the g-and-k result above. In one dimension the effective sample size is not just small but \emph{stable}: $237$--$297$ at every checkpoint of every replicate, and accuracy is correspondingly steady, with no relationship between the two (log--log correlation $0.04$). In four dimensions it is small \emph{and} erratic --- across the five g-and-k replicates it ranges from $7$ to $2987$ at the same points in wall-clock time, moving by two orders of magnitude between neighbouring checkpoints of a single run --- and accuracy follows it closely: the log--log correlation between effective sample size and $W_1$ to the reference posterior is $-0.80$. Every excursion above $W_1 = 0.3$ in that panel is a checkpoint where the effective sample had collapsed below $50$. So the higher-dimensional deficit is not a bias in the estimator but intermittent weight degeneracy in the cumulative-mixture denominator, and it is the same effect the calibration study meets as an under-covering top-$k$ report and an over-covering full-history one (\S\ref{sec:results-validity}), seen here as variance rather than as coverage.

The condition for the $\mathrm{ESS} \approx 3k$ ceiling is that the bandwidth actually tightens, which is to say that the target is identified --- and the Cellular Potts run is the instructive counter-case. There the reported posterior retains an ESS \emph{fraction} of $0.21$ over $33{,}748$ evaluations, i.e.\ ${\approx}7{,}200$ effective particles against the Gaussian run's ${\approx}280$, because the parameters confound (\S\ref{sec:results-cpm}), the bandwidth stalls, and almost every particle keeps weight. The two cases are the same mechanism read at its two ends, and neither is good news on its own: an identified target concentrates the effective support onto a few multiples of $k$, and a confounded one keeps a large effective sample of an uninformative posterior. This is the same object as the stabilization rate of \S\ref{sec:theory-r}, seen from the estimator's side rather than the bandwidth's: what governs it is the identifiability of the target, not the sampler.

We think this is the single most useful thing to know about the method, and it reframes three results we had reported separately. It is why reporting over the full history rather than the top-$k$ archive helps calibration but only up to a point, and why an intermediate reported support of $M\approx400$--$800$ was where coverage crossed nominal: that is $4$--$8k$, about where the effective support already sits. It is consistent with the measured failure of the CLT's rate condition for the reported configuration (\S\ref{sec:theory-r}): an asymptotic variance that shrinks at rate $n$ is not available from an effective sample that does not grow with $n$. And it says plainly where the method's remaining headroom is --- in $k$ and in the bandwidth rule, not in more workers. The throughput results of \S\ref{sec:results-hpc} stand on their own as systems results; what they buy in posterior accuracy, on a benchmark this well specified, is bounded by an effective support that more simulation alone does not enlarge.

\begin{figure}[t]
\centering
\includegraphics[width=\linewidth]{figures/fig_posterior_recovery.pdf}
\caption{Posterior quality (Wasserstein distance to the truth/reference, lower is better) versus wall-clock time on g-and-k (left), Lotka--Volterra (middle), and Cellular Potts (right), asynchronous vs.\ synchronous; shaded bands are inter-quartile ranges over five replicates. On g-and-k the asynchronous method reaches a slightly lower Wasserstein distance within the budget, on Lotka--Volterra the two are comparable, and on Cellular Potts the asynchronous method is modestly but consistently behind. Note that on the weakly identified Cellular Potts target the distance-to-reference \emph{increases} with compute for both methods, so the metric penalizes the method that concentrates fastest (\S\ref{sec:results-validity}).}\label{fig:posterior-recovery}
\end{figure}

\begin{figure}[t]
\centering
\includegraphics[width=0.55\linewidth]{figures/fig_cpm_recovery_diff.pdf}
\caption{Cellular Potts posterior-quality \emph{difference} (asynchronous minus synchronous Wasserstein to the reference) versus wall-clock time, with a zero reference line and an inter-quartile envelope over five replicates (the per-method Cellular Potts curves of Fig.~\ref{fig:posterior-recovery}c, differenced). The difference is positive --- asynchronous behind --- across $87.5\%$ of the budget, ending at ${+}0.043$ with an envelope of $[0.019,0.052]$ that excludes zero; it straddles zero only over the first ${\approx}510$\,s. The gap is nonetheless small in absolute terms, never exceeding ${\approx}0.06$ against an absolute Wasserstein of ${\approx}0.4$, and it is measured against a reference point on a weakly identified target whose distance-to-reference grows as either method concentrates (\S\ref{sec:results-validity}, \S\ref{sec:results-cpm}).}\label{fig:cpm-diff}
\end{figure}

\begin{figure}[t]
\centering
\includegraphics[width=\linewidth]{figures/fig_reported_recovery.pdf}
\caption{What the \emph{reported} estimator recovers, and with how many effective particles. Each method is scored on the posterior it actually reports --- for the asynchronous method the retroactive AMIS estimator replayed over the prefix completed by each checkpoint, for the synchronous baseline its last completed generation, for rejection ABC its accepted set --- against a reference posterior rather than a point mass at the truth. \textbf{Top:} $W_1$ to the reference, median and inter-quartile range over five replicates. \textbf{Bottom:} the effective sample size of that same reported posterior, with the evaluated-particle count (dotted) on the same axis. The gap between the two lines is the point: the history grows by more than an order of magnitude and the effective sample does not follow it.}\label{fig:reported-recovery}
\end{figure}

\subsection{Realistic simulator: Cellular Potts (RQ3)}\label{sec:results-cpm}
The Cellular Potts benchmark runs end-to-end with the real \texttt{cellsInSilico} simulator, a realistic, expensive, heterogeneous-runtime workload and the regime the method targets.

\paragraph{Strong scaling.} Because each Cellular Potts evaluation carries real cost, the per-arrival coordination that capped Lotka--Volterra is amortized, and the asynchronous method scales almost linearly all the way to eight nodes (Fig.~\ref{fig:scaling}, right): median throughput rises from $0.2$ to $72.0$ simulations/s over $1$ to $384$ workers and, unlike the near-instantaneous Lotka--Volterra case, keeps scaling \emph{across} the single-node$\to$multi-node boundary, each node-doubling roughly doubling throughput ($9.3\to18.7\to37.8\to72.0$ sims/s at $48\to96\to192\to384$ workers), for $\approx 97\%$ parallel efficiency relative to the single-node rate. The synchronous baseline, given a population equal to the worker count so it too can occupy every core, scales only \emph{sub-linearly} over the same range ($3.8\to4.5\to7.8\to14.2$ sims/s at $48\to96\to192\to384$ workers, ${\approx}50\%$ efficiency), so at $384$ workers the asynchronous method sustains ${\approx}5\times$ its throughput, a gap that \emph{widens} with scale. This is the regime the method is built for: an expensive, heterogeneous simulator where the generation barrier, not per-arrival overhead, is the binding cost. Worker utilization makes the mechanism explicit (Fig.~\ref{fig:cpm-util}): the asynchronous method keeps workers ${\approx}98.9$--$99.8\%$ busy across all scales, whereas the synchronous baseline, even with every core filled within a generation, then waits for the slowest simulation at the barrier, so its utilization falls to ${\approx}24$--$29\%$ ($41\%\to24\%\to26\%\to29\%$ from $48$ to $384$ workers, i.e.\ ${\approx}70$--$76\%$ idle at every multi-node scale). This barrier idle, not per-arrival overhead, is what holds the synchronous throughput to sub-linear scaling while the asynchronous method tracks the ideal. (We report medians throughout.)

\begin{figure}[t]
\centering
\includegraphics[width=0.52\linewidth]{figures/fig_cpm_util.pdf}
\caption{Cellular Potts worker utilization (fraction of worker wall-clock spent inside the simulator) versus worker count ($1800$\,s budget; error bars are $\pm1$ s.d.); the synchronous baseline uses a population equal to the worker count (Fig.~\ref{fig:scaling}, right). The asynchronous method stays ${\approx}98.9$--$99.8\%$ utilized at every scale, while the synchronous baseline, even filling every core within a generation, then idles at the barrier waiting for the slowest simulation, so its utilization falls to ${\approx}24$--$29\%$ ($41\%\to24\%\to26\%\to29\%$ at $48/96/192/384$ workers, i.e.\ ${\approx}70$--$76\%$ idle at every multi-node scale). This barrier idle is the measured mechanism behind the synchronous method's sub-linear scaling in Fig.~\ref{fig:scaling} (right).}\label{fig:cpm-util}
\end{figure}

\paragraph{The barrierized twin at this scale.} The twin comparison of \S\ref{sec:results-hpc} isolates synchronization on the two synthetic benchmarks; here we run it on the real simulator, at the same four worker counts and the same $k{=}100$, simulation-limited at $40$ evaluations per worker (Table~\ref{tab:twin-cpm}). Two things come out of it, and the second is the more interesting.

First, the barrier alone costs a factor that is remarkably \emph{stable} across scale: $1.9\times$ at one, two and four nodes, rising to $2.8\times$ at eight. The external pyABC baseline's penalty over the same range grows much faster, from $2.4\times$ to $5.1\times$. The twin is therefore consistently the faster of the two synchronous arms, and the distance between them widens with scale --- which is the honest answer to the question the twin was built to ask, and it is not the flattering one. On this workload the barrier is a \emph{major but not exclusive} cause of the cross-algorithm gap: it accounts for roughly half to three-quarters of it in log terms, and the residual, which grows with worker count, has to come from pyABC's other differences (its probabilistic-rejection step, its classical population weighting, and its own coordination costs) rather than from synchronization. We therefore do \emph{not} claim the Cellular Potts gap is attributable to the barrier alone --- only that the barrier is the single largest identified component of it. This is also the opposite ordering to the straggler and heterogeneity benchmarks, where the twin ran far \emph{slower} than pyABC because pyABC's latecomer discarding rescues it from waiting on the slowest draw; with no injected pathology, that rule has little to exploit and our more efficient propagator keeps the twin ahead.

Second --- and this is the more useful finding --- the barrier does not merely reduce throughput, it makes throughput \emph{unpredictable}. The asynchronous arm's five replicates agree to within $1.3\%$ at every scale (coefficient of variation $0.6$--$1.3\%$), while the twin's spread across the \emph{same} seeds is $25$--$54\%$: at $48$ workers its replicates run at $1.7$, $2.1$, $4.9$, $5.9$ and $6.0$ simulations/s, a factor of $3.5$ between fastest and slowest. The mechanism is the one the utilization measurement already shows: a generation ends only when its slowest evaluation does, so the whole job inherits the cost of the most expensive parameter draw in each generation, and on this benchmark simulation cost varies strongly with the parameters. Asynchrony removes that coupling --- a slow draw delays only its own worker --- which is why the asynchronous arm's throughput is reproducible to within a percent while its barrierized twin, running identical proposals from identical seeds, varies by more than threefold. For a practitioner this is arguably the more consequential of the two effects: a median slowdown can be budgeted for, a factor-of-three variance in time-to-result cannot.

One note on the twin as an instrument. It intermittently deadlocks at $\ge192$ ranks --- we observed five hangs while collecting this table, one in the MPI teardown and the others mid-run --- so individual replicates had to be retried. Every combination did complete, and all twenty entries here are full runs; but the twin is a diagnostic built for this comparison rather than a robust implementation, and at these rank counts it needs supervision.

\begin{table}[t]
\caption{Barrierized twin on the Cellular Potts workload: median throughput (simulations/s) over five replicates, $k{=}100$, simulation-limited at $40$ evaluations per worker. CV is the coefficient of variation across replicates. The barrier's own cost is nearly flat from one to four nodes ($1.9\times$) and rises to $2.8\times$ at eight, while the external pyABC baseline's penalty grows from $2.4\times$ to $5.1\times$ over the same range --- so the barrier is the single largest identified component of the cross-algorithm gap but not the whole of it, and the residual grows with scale. The second row-pair is the more useful finding: the twin's replicate spread is $25$--$54\%$ against the asynchronous arm's $0.6$--$1.3\%$ on the same seeds, because a generation inherits the cost of its most expensive parameter draw. Synchronization makes throughput unpredictable as well as lower. The pyABC baseline has three replicates at $192$ and $384$ workers (its fair population$=W$ shards); all twin and asynchronous entries have five.}\label{tab:twin-cpm}
\centering
\begin{tabular}{@{}lcccc@{}}
\toprule
Workers ($=$ cores) & 48 (1 node) & 96 (2) & 192 (4) & 384 (8) \\
\midrule
Asynchronous (ours)     & 9.33 & 18.67 & 37.75 & 71.98 \\
\quad CV                & $0.9\%$ & $1.3\%$ & $0.6\%$ & $1.3\%$ \\
Barrierized twin        & 4.87 & 10.00 & 20.24 & 25.54 \\
\quad CV                & $50.7\%$ & $53.8\%$ & $24.8\%$ & $49.8\%$ \\
Synchronous pyABC baseline & 3.85 & 4.53 & 7.85 & 14.17 \\
\midrule
Ratio, twin             & $1.9\times$ & $1.9\times$ & $1.9\times$ & $2.8\times$ \\
Ratio, pyABC baseline   & $2.4\times$ & $4.1\times$ & $4.8\times$ & $5.1\times$ \\
\botrule
\end{tabular}
\end{table}

\paragraph{Posterior quality.} On the inference side the asynchronous method produces marginals of similar shape to the synchronous baseline (Fig.~\ref{fig:cpm-posterior}), though it trails modestly on distance-to-reference (\S\ref{sec:results-validity}): the cell-motility parameter concentrates near its reference value, while the division rate is not merely unresolved but mislocated --- every method, the baseline and rejection ABC included, places its mass away from the reference (\S\ref{sec:results-validity}). That is a property of the configuration rather than an incidental finding: this benchmark was built to supply a realistic heterogeneous-runtime workload for the scaling and throughput measurements, and its summaries, discrepancy metric and parameter ranges were never selected for identifiability, with the two parameters confounding. Quality comparisons on it should be read as evidence that no method is advantaged by the setup, not as a measurement of what any of them can achieve. This weak identification is a structural property of the inference problem, not of asynchrony: at this scale the division-rate and motility parameters are near-degenerate (both chiefly modulate the cell count that the summaries resolve), so neither method separates them and both improve only modestly on rejection ABC. The matched comparison therefore isolates the systems advantage as largely \emph{orthogonal} to identifiability: both methods are limited chiefly by the problem. On the summary metric the asynchronous method is modestly behind over the full budget (${+}0.043$ in Wasserstein-to-reference against a scale of ${\approx}0.4$; \S\ref{sec:results-validity}), so we report its throughput advantage as bought at a small quality cost on this benchmark rather than at none --- with the caveat that the metric rewards dispersion on a weakly identified target, and that the asynchronous method, concentrating fastest, is the one it penalizes most. Establishing the same quality with a synchronous matched baseline at the still larger scales the method targets would itself be prohibitively wasteful, precisely the cost the barrier-free design removes.

\begin{figure}[t]
\centering
\includegraphics[width=0.62\linewidth]{figures/fig_cpm_corner.pdf}
\caption{Cellular Potts posterior (division rate, motility) on the real \texttt{cellsInSilico} simulator: asynchronous (blue), synchronous baseline (red), and rejection ABC (green), with dashed lines at the reference values. The asynchronous posterior is drawn from its AMIS estimator ($500$ resampled draws); the synchronous baseline and rejection ABC show their final populations ($500$ each). The asynchronous method and the synchronous baseline both place their motility mass at low values, bracketing the motility reference (which sits at the $87$th and $91$st percentiles of their marginals). Neither recovers the division rate: its reference lies below every one of the $500$ draws for both, and rejection ABC puts it at its own $5$th percentile --- the summaries and parameter ranges of this benchmark were never selected for identifiability (\S\ref{sec:results-cpm}), so the panel bounds no method's attainable accuracy. Both do concentrate more tightly than rejection ABC.}\label{fig:cpm-posterior}
\end{figure}

\subsection{Ablation}\label{sec:results-ablation}
Reverting the smooth kernel to a hard indicator slightly raises the Wasserstein-to-truth relative to the full method (though within its confidence interval), while removing the AMIS cumulative-mixture reweighting is approximately neutral on this uniform-runtime analytic target (Fig.~\ref{fig:ablation}). On the Gaussian-mean benchmark the full method attains a final Wasserstein-to-truth of $0.073$ (CI $[0.067,0.079]$); the hard-indicator kernel raises it to $0.076$ (inside that interval), whereas dropping AMIS leaves it essentially unchanged ($0.072$; Fig.~\ref{fig:ablation}a), with the remaining hyperparameter variants scattered around the full-method reference and degrading only toward the largest perturbation scales ($0.084$--$0.087$). That AMIS is neutral here is consistent with the coupling study (\S\ref{sec:results-hpc}), where removing AMIS likewise left the posterior essentially unchanged: with no runtime--parameter coupling on this uniform-runtime target, there is no over-representation for the reweighting to reshape. The AMIS-isolation trajectory (Fig.~\ref{fig:ablation}b) accordingly shows the full method and the no-AMIS variant tracking each other throughout the converged regime (both $\approx 0.072$--$0.073$ after $40$\,s).

\begin{figure}[t]
\centering
\includegraphics[width=\linewidth]{figures/fig_ablation.pdf}
\caption{Ablation on the Gaussian-mean benchmark (analytic posterior; $48$ workers, $300$\,s budget, $k=100$, five replicates, asynchronous method). (a) Final Wasserstein-to-truth for the full method (blue, solid), the two ingredient removals (vermillion, hatched: no AMIS, hard kernel), and hyperparameter variants (grey, cross-hatched); colour and hatch encode the class redundantly, error bars are $95\%$ CIs, and the dotted line marks the full-method mean. (b) AMIS isolation: full method vs.\ no-AMIS on a shared wall-clock grid (bands are $\pm1$ s.d.\ over replicates), zoomed to the converged regime; on this uniform-runtime analytic target the two variants track each other.}\label{fig:ablation}
\end{figure}

\subsection{Sensitivity}\label{sec:results-sensitivity}
Across the full hyperparameter grid \emph{on the Gaussian-mean benchmark} (the analytic-posterior sanity model) the asynchronous method's posterior quality is stable, with no region of catastrophic failure (Fig.~\ref{fig:sensitivity}). The Wasserstein distance to the true parameter stays within $\approx0.09$--$0.15$ across the grid (each cell averaged over five replicates and the two smooth kernels); the archive size $k$ matters only marginally \emph{for point quality at this dimension} (its effect on calibration in higher dimensions is a separate and much larger story, Table~\ref{tab:ks-sweep}), and quality degrades only mildly toward the largest perturbation scales and initial tolerances. The tolerance scheduler is inert under a smooth kernel --- the asynchronous selector chooses $\epsilon$ by ESS-retention bisection regardless of the scheduler type --- so we fixed it to the acceptance-rate rule rather than sweeping a variable that cannot change the smooth-kernel result. We therefore did not find delicate tuning necessary to reach the quality reported above \emph{on this benchmark}. We did not repeat the full grid on the costlier non-Gaussian benchmarks, so this robustness statement is scoped to the analytic-posterior model and to the qualitative trends observed there (scheduler-insensitive; mild, monotone sensitivity to perturbation scale and initial tolerance) rather than asserted as a general guarantee. Two further scope limits matter. This grid measures \emph{point quality} (Wasserstein to the truth) in \emph{one} dimension; it says nothing about calibration, and the archive size that looks inconsequential here is decisive for coverage in higher dimensions --- the dimension-resolved sweep of \S\ref{sec:results-validity} (Table~\ref{tab:ks-sweep}) finds a genuinely catastrophic region ($d=8$, $k=50$: coverage $0.29$ at the $0.50$ level) that a one-dimensional quality grid cannot see. ``No failure region'' should therefore be read as a statement about one-dimensional point quality only.

\begin{figure}[t]
\centering
\includegraphics[width=\linewidth]{figures/fig_sensitivity_heatmap.pdf}
\caption{Hyperparameter sensitivity of the asynchronous method, measured by posterior quality (Wasserstein distance to the true $\mu$, lower is better; averaged over replicates and smoothing kernel). Each panel sweeps perturbation scale ($x$) against the initial-tolerance multiplier ($y$); panels are faceted by archive size $k$. Quality is stable across the grid (Wasserstein $\approx0.09$--$0.15$) with no failure region; the tolerance scheduler is inert under the smooth kernel and is held fixed, while larger perturbation scales degrade quality only mildly.}\label{fig:sensitivity}
\end{figure}

